\chapter{Fonctions de deux variables}\label{ch:b1:multivar}

L'année s'achève par une première incursion en dimension
supérieure~: les fonctions $f(x, y)$ de deux variables réelles.
Tout se généralise --- limites, \omterm{def:b1:continuity:continuous}{continuité}, \omterm{def:b1:derivative:def}{dérivées}, extrema ---
mais chaque notion prend un tour nouveau~: une limite peut être
approchée selon toutes les directions à la fois, les \omterm{def:b1:derivative:def}{dérivées} se
scindent en \omterm{def:b1:multivar:partial}{dérivées partielles}, et le \omterm{def:b1:multivar:partial}{gradient} indique la
direction de la montée. La théorie complète (différentielles,
$\R^n$ général, sous-variétés) relève de la deuxième année~; ici
nous fixons le vocabulaire et les premiers théorèmes honnêtes.

\section{Le plan $\R^2$~; continuité}

\begin{definition}\label{def:b1:multivar:topology}
Sur $\R^2$, on utilise la norme euclidienne $\norm{(x,y)} = \sqrt{x^2 + y^2}$
(\cref{ch:b1:euclid}). Les boules \omterm{def:b1:topology:open}{ouvertes}, les \omterm{def:b1:topology:open}{voisinages}, les
\emph{parties \omterm{def:b1:topology:open}{ouvertes}} de $\R^2$ se définissent exactement comme
au \cref{ch:b1:topology}, avec des boules à la place des
\omterm{prop:b1:reals:intervals}{intervalles}. Une fonction $f \colon U \to \R$ ($U \subseteq \R^2$
\omterm{def:b1:topology:open}{ouvert}) est \emph{\omterm{def:b1:continuity:continuous}{continue} en} $a \in U$ lorsque
\[
\forall\varepsilon > 0,\ \exists\delta > 0, \quad
\norm{X - a} \leq \delta \implies \abs{f(X) - f(a)} \leq
\varepsilon,
\]
avec la même caractérisation séquentielle qu'en une variable.
Sommes, produits, quotients et composées avec des fonctions
\omterm{def:b1:continuity:continuous}{continues} d'une variable préservent la \omterm{def:b1:continuity:continuous}{continuité}~; les
\omterm{def:b1:logic:map}{applications} \omterm{prop:b1:vspaces:coordinates}{coordonnées} sont \omterm{def:b1:continuity:continuous}{continues}, donc aussi les \omterm{def:b1:poly:def}{polynômes}
en $(x,y)$.
\end{definition}

\begin{example}[La majoration polaire, la bonne façon de prouver une limite]
\label{ex:b1:multivar:polarbound}
Montrons que $f(x, y) = \dfrac{x^2y^2}{x^2 + y^2}$ (avec $f(0,0) =
0$) est \omterm{def:b1:continuity:continuous}{continue} à l'origine. En \omterm{prop:b1:vspaces:coordinates}{coordonnées} polaires $x =
\rho\cos\theta$, $y = \rho\sin\theta$~:
\[
\abs{f} = \frac{\rho^4\cos^2\theta\sin^2\theta}{\rho^2}
= \rho^2\,(\cos\theta\sin\theta)^2 \leq \frac{\rho^2}{4}
\xrightarrow[\rho \to 0]{} 0 ,
\]
majoration \emph{indépendante de $\theta$}~: quelle que soit la
direction d'approche, les valeurs sont écrasées vers $0$. Cette
uniformité en $\theta$ est tout l'enjeu --- une majoration comme
$\abs g = \abs{\cos\theta\sin\theta}$ (plus de $\rho$) ne prouve
rien, et de fait ce $g$ est le piège radial discontinu de
l'exemple suivant.
\end{example}

\begin{example}[Le piège radial]\label{ex:b1:multivar:trap}
Soit $f(x, y) = \dfrac{xy}{x^2 + y^2}$ pour $(x,y) \neq (0,0)$,
$f(0, 0) = 0$. Le long de chaque axe, $f = 0 \to 0$~; mais le long
de la diagonale $y = x$, $f(x, x) = \frac12 \not\to 0$. \emph{Pas
de limite à l'origine}~: approcher selon toutes les droites, et
même trouver la même limite selon chacune, ne suffit pas (ici les
limites selon les droites diffèrent~; il y a pire, des fonctions
qui donnent la même limite selon toutes les droites et échouent le
long d'une parabole, \cref{exo:b1:multivar:3}). La \omterm{def:b1:continuity:continuous}{continuité} par
rapport à chaque variable séparément n'entraîne pas la \omterm{def:b1:continuity:continuous}{continuité}.
\end{example}

\section{Dérivées partielles}

\begin{definition}\label{def:b1:multivar:partial}
Les \emph{dérivées partielles}\index{dérivée partielle} de $f$ en
$(a, b)$ sont les \omterm{def:b1:derivative:def}{dérivées} à une variable le long des axes~:
\[
\frac{\partial f}{\partial x}(a,b)
= \lim_{h \to 0} \frac{f(a + h, b) - f(a,b)}{h},
\qquad
\frac{\partial f}{\partial y}(a,b)
= \lim_{k \to 0} \frac{f(a, b + k) - f(a,b)}{k}.
\]
$f$ est de classe $C^1$ sur $U$ lorsque les deux existent et sont
\omterm{def:b1:continuity:continuous}{continues} sur $U$. Le \emph{gradient}\index{gradient} est $\nabla
f(a,b) = \bigl(\frac{\partial f}{\partial x},\, \frac{\partial
f}{\partial y}\bigr)(a,b)$.
\end{definition}

\begin{theorem}[$C^1$ entraîne un plan tangent]\label{thm:b1:multivar:c1}
Soit $f$ de classe $C^1$ sur $U$ et $(a,b) \in U$. Alors, quand
$(h, k) \to (0,0)$~:
\[
f(a + h, b + k) = f(a, b)
+ h\,\frac{\partial f}{\partial x}(a,b)
+ k\,\frac{\partial f}{\partial y}(a,b)
+ o\bigl(\norm{(h,k)}\bigr) .
\]
En particulier $f$ est \omterm{def:b1:continuity:continuous}{continue}, et le graphe $z = f(x,y)$ admet en
chaque point le plan tangent que la formule fait lire.
\end{theorem}

\begin{example}[L'approximation linéaire à l'œuvre]
\label{ex:b1:multivar:linapprox}
Estimons $f(1.02,\ 0.99)$ pour $f(x, y) = x^3y^2$. En $(1, 1)$~:
$f = 1$, $\frac{\partial f}{\partial x} = 3x^2y^2 = 3$,
$\frac{\partial f}{\partial y} = 2x^3y = 2$, donc
le \cref{thm:b1:multivar:c1} donne
\[
f(1.02,\ 0.99) \approx 1 + 3\,(0.02) + 2\,(-0.01) = 1.04 ,
\]
contre la valeur exacte $1.02^3 \times 0.99^2 = 1.04006\dots$ ---
l'erreur est du second ordre en les accroissements, comme le
promet le $o(\norm{(h,k)})$. Le plan tangent au graphe en
$(1, 1, 1)$ est $z = 1 + 3(x - 1) + 2(y - 1)$, l'équation
implicitement utilisée dans l'estimation.
\end{example}

\begin{proof}
Déplaçons une coordonnée à la fois~:
\[
f(a+h, b+k) - f(a,b)
= \bigl[f(a+h, b+k) - f(a, b+k)\bigr] + \bigl[f(a, b+k) -
f(a,b)\bigr].
\]
Par le \omterm{thm:b1:derivative:mvt}{théorème des accroissements finis} à une variable
(\cref{thm:b1:derivative:mvt}), le premier crochet vaut $h\,
\frac{\partial f}{\partial x}(a + \theta h,\, b + k)$ pour un
certain $\theta \in \intoo{0}{1}$, et le second $k\,\frac{\partial
f}{\partial y}(a,\, b + \theta' k)$. La \omterm{def:b1:continuity:continuous}{continuité} des \omterm{def:b1:multivar:partial}{dérivées
partielles} en $(a,b)$ permet d'écrire chacune comme (valeur en
$(a,b)$) $+$ (erreur $\to 0$)~; l'erreur totale est
$h\,\varepsilon_1 + k\,\varepsilon_2 = o(\norm{(h,k)})$ puisque
$\abs h, \abs k \leq \norm{(h,k)}$.
\end{proof}

\begin{theorem}[Règle de la chaîne]\label{thm:b1:multivar:chain}
Soit $f$ de classe $C^1$ sur $U$ et $t \mapsto (x(t), y(t))$ de
classe $C^1$ d'un \omterm{prop:b1:reals:intervals}{intervalle} dans $U$. Alors $g(t) =
f\bigl(x(t), y(t)\bigr)$ est de classe
$C^1$, et\index{règle de la chaîne!deux variables}
\[
g'(t) = x'(t)\,\frac{\partial f}{\partial x}\bigl(x(t),y(t)\bigr)
+ y'(t)\,\frac{\partial f}{\partial y}\bigl(x(t),y(t)\bigr)
= \bigl\langle \nabla f,\ (x', y')\bigr\rangle .
\]
\end{theorem}

\begin{proof}
Appliquons le \cref{thm:b1:multivar:c1} en $(x(t), y(t))$ avec $(h, k) =
(x(t+s) - x(t),\, y(t+s) - y(t))$~: quand $s \to 0$, la
\omterm{def:b1:derivative:def}{dérivabilité} à une variable donne $h = s\,x'(t) + o(s)$ et $k =
s\,y'(t) + o(s)$, donc $\norm{(h, k)} = O(s)$ et
\[
g(t+s) - g(t)
= h\,\frac{\partial f}{\partial x} + k\,\frac{\partial
f}{\partial y} + o\bigl(\norm{(h,k)}\bigr)
= s\,\Bigl(x'\,\frac{\partial f}{\partial x} +
y'\,\frac{\partial f}{\partial y}\Bigr) + o(s) ,
\]
l'erreur finale absorbant à la fois les $o(s)$ de $h$ et de $k$
(multipliés par les valeurs fixes des \omterm{def:b1:multivar:partial}{dérivées partielles}) et le
$o(O(s))$ de l'estimation du plan tangent. On divise par $s$ et
l'on fait tendre $s \to 0$. La \omterm{def:b1:continuity:continuous}{continuité} de $g'$ résulte de celle
de tous les ingrédients.
\end{proof}

\begin{example}[La règle de la chaîne, vérifiée des deux façons]
\label{ex:b1:multivar:chainrun}
Soit $f(x, y) = x^2 y$ et $g(t) = f(t, t^2)$. Directement~: $g(t) =
t^2\cdot t^2 = t^4$, donc $g'(t) = 4t^3$. Par la \omterm{thm:b1:derivative:operations}{règle de la
chaîne}~: $\frac{\partial f}{\partial x} = 2xy$ et $\frac{\partial
f}{\partial y} = x^2$, évaluées le long de la courbe $(t, t^2)$~:
\[
g'(t) = 1\cdot 2t\cdot t^2 + 2t\cdot t^2 = 2t^3 + 2t^3 = 4t^3 .
\]
Les deux calculs concordent, et la décomposition a un sens~: $2t^3$
de la croissance vient du déplacement \emph{vers la droite} à
travers la pente en $x$, $2t^3$ du déplacement \emph{vers le haut}
à travers la pente en $y$. Sur les courbes où $g$ n'a pas de forme
close, seul le second calcul survit --- c'est là tout l'intérêt du
théorème.
\end{example}

\begin{remark}[Lire le gradient]
Le long d'une direction unitaire $u$, la \omterm{thm:b1:derivative:operations}{règle de la chaîne}
appliquée à $t \mapsto f(a + tu)$ donne la \emph{\omterm{def:b1:derivative:def}{dérivée}
directionnelle} $\langle \nabla f(a), u\rangle$~: maximale quand
$u$ pointe dans la direction de $\nabla f(a)$ (par
Cauchy--Schwarz, \cref{thm:b1:euclid:cs}). Le \omterm{def:b1:multivar:partial}{gradient} est la
direction de plus forte pente, et il est \omterm{def:b1:euclid:orthogonal}{orthogonal} aux lignes de
niveau $\{f = c\}$ (dériver $f$ le long d'une courbe tracée dans
une ligne de niveau~: la \omterm{thm:b1:derivative:operations}{règle de la chaîne} donne $\langle\nabla
f,\, \text{tangente}\rangle = 0$).
\end{remark}

\begin{example}[Lignes de niveau et gradients, sur une même fonction]
\label{ex:b1:multivar:levelcurves}
Prenons $f(x, y) = x^2 - y^2$. Ses lignes de niveau~: $\{f = c\}$
est une hyperbole \omterm{def:b1:topology:open}{ouverte} vers la gauche et la droite pour $c >
0$, vers le haut et le bas pour $c < 0$, et le couple de droites
sécantes $y = \pm x$ pour $c = 0$ --- la carte des courbes de
niveau d'un col de montagne, avec le \omterm{met:b1:multivar:monge}{point col} à l'origine où les
deux droites de niveau nul se croisent. \omterm{def:b1:multivar:partial}{Gradient}~: $\nabla
f = (2x, -2y)$. Au point $(2, 1)$ (sur le niveau $c = 3$)~:
$\nabla f = (4, -2)$, tandis que le vecteur tangent à la ligne de
niveau, paramétrée près de ce point par $\bigl(t,
\sqrt{t^2 - 3}\bigr)$, vaut $\bigl(1, \frac{t}{\sqrt{t^2 -
3}}\bigr) = (1, 2)$ en $t = 2$ --- et de fait
\[
\langle (4, -2),\ (1, 2)\rangle = 4 - 4 = 0 :
\]
\omterm{def:b1:multivar:partial}{gradient} perpendiculaire à la courbe de niveau, pointant vers les
valeurs plus grandes de $f$ (ici~: en s'éloignant de l'axe des
$y$). Deux lectures de plus~: le \omterm{def:b1:multivar:partial}{gradient} s'annule exactement au
\omterm{met:b1:multivar:monge}{point col}, là où la carte des niveaux se pince~; et la
\emph{droite} tangente à la ligne de niveau en $(2,1)$ est
$4(x - 2) - 2(y - 1) = 0$, c'est-à-dire $2x - y = 3$ ---
l'équation «~$\langle \nabla f,\ M - M_0\rangle = 0$~» qui
généralise la tangente à l'ellipse de l'\cref{exo:b1:curves:11}.
\end{example}

\begin{theorem}[Schwarz]\label{thm:b1:multivar:schwarz}
Si $f$ est de classe $C^2$ (les \omterm{def:b1:multivar:partial}{dérivées partielles} des \omterm{def:b1:multivar:partial}{dérivées
partielles} existent et sont \omterm{def:b1:continuity:continuous}{continues}),
alors\index{théorème de Schwarz}
\[
\frac{\partial^2 f}{\partial x\,\partial y}
= \frac{\partial^2 f}{\partial y\,\partial x} .
\]
\end{theorem}
\admitted

\section{Extrema locaux}

\begin{method}[Étude des extrema, en bon ordre]
\label{met:b1:multivar:extremamethod}
\begin{enumerate}
  \item \emph{Résoudre $\nabla f = 0$ complètement.} Factoriser
        chaque \omterm{def:b1:multivar:partial}{dérivée partielle} dès que possible (des produits
        de facteurs \omterm{def:b1:linmaps:def}{linéaires} scindent le système en cas
        transparents, comme dans
        l'\cref{ex:b1:multivar:fourpoints} ci-dessous)~; un cas
        oublié est un \omterm{thm:b1:multivar:critical}{point critique} oublié.
  \item \emph{Classer chaque point} à l'aide des données de
        Monge $r, s, t$ --- recalculées \emph{en chaque point},
        jamais une fois pour toutes.
  \item \emph{Si $rt - s^2 = 0$}, examiner $f$ directement le
        long de courbes bien choisies passant par le point (les
        droites d'abord, puis les paraboles), en cherchant soit
        deux signes (pas d'extremum), soit un signe constant avec
        un argument couvrant toutes les directions.
  \item \emph{Prendre du recul pour la vue d'\omterm{def:b1:logic:sets}{ensemble}}~:
        examiner le comportement à l'infini (un minimum local
        peut n'être pas global), et si le domaine n'est pas
        \omterm{def:b1:topology:open}{ouvert}, traiter sa \omterm{def:b1:topology:closure}{frontière} à part
        (\cref{exo:b1:multivar:12}) --- le théorème des \omterm{thm:b1:multivar:critical}{points
        critiques} ne voit que les points \omterm{def:b1:topology:closure}{intérieurs}.
\end{enumerate}
\end{method}

\begin{theorem}[Points critiques]\label{thm:b1:multivar:critical}
Si $f$ (de classe $C^1$ sur l'\omterm{def:b1:topology:open}{ouvert} $U$) admet un extremum local
en $(a,b) \in U$, alors $\nabla f(a,b) = (0,0)$~: le point est
\emph{critique}\index{point critique}.
\end{theorem}

\begin{proof}
Les fonctions d'une variable $x \mapsto f(x, b)$ et $y \mapsto f(a,
y)$ admettent des extrema locaux \omterm{def:b1:topology:closure}{intérieurs} en $a$, resp.\ $b$~:
la \cref{prop:b1:derivative:fermat} annule les deux \omterm{def:b1:multivar:partial}{dérivées
partielles}.
\end{proof}

\begin{method}[Test du second ordre (notations de Monge)]\label{met:b1:multivar:monge}
En un \omterm{thm:b1:multivar:critical}{point critique} d'une fonction de classe $C^2$, posons
\[
r = \frac{\partial^2 f}{\partial x^2},
\qquad
s = \frac{\partial^2 f}{\partial x \partial y},
\qquad
t = \frac{\partial^2 f}{\partial y^2}
\qquad (\text{valeurs au point}).
\]
\begin{itemize}
  \item Si $rt - s^2 > 0$~: extremum local --- minimum pour $r > 0$,
        maximum pour $r < 0$~;
  \item si $rt - s^2 < 0$~: pas d'extremum (un \emph{point
        col}\index{point col})~;
  \item si $rt - s^2 = 0$~: le test est muet~; examiner
        directement.
\end{itemize}
(La justification --- une formule de Taylor--Young à l'ordre $2$ et
l'étude du signe de la forme quadratique $r h^2 + 2shk + tk^2$ ---
est menée en deuxième année~; le test est utilisé ici comme outil
de travail.)
\end{method}

\begin{example}\label{ex:b1:multivar:extrema}
$f(x,y) = x^3 + y^3 - 3xy$. \omterm{thm:b1:multivar:critical}{Points critiques}~: $\nabla f = (3x^2 -
3y,\; 3y^2 - 3x) = 0$ donne $y = x^2$ et $x = y^2$, donc $x = x^4$~:
$x \in \{0, 1\}$~: points $(0,0)$ et $(1,1)$.

\omterm{def:b1:derivative:def}{Dérivées} secondes~: $r = 6x$, $s = -3$, $t = 6y$. En $(0,0)$~: $rt -
s^2 = -9 < 0$~: \omterm{met:b1:multivar:monge}{point col}. En $(1,1)$~: $rt - s^2 = 36 - 9 > 0$, $r = 6
> 0$~: minimum local, $f(1,1) = -1$. (Non global~: $f(x, 0) = x^3 \to
-\infty$.)
\end{example}

\begin{example}[Une étude à quatre points, en entier]
\label{ex:b1:multivar:fourpoints}
$f(x, y) = xy\,(3 - x - y) = 3xy - x^2y - xy^2$. \omterm{def:b1:multivar:partial}{Gradient}~:
\[
\frac{\partial f}{\partial x} = y\,(3 - 2x - y),
\qquad
\frac{\partial f}{\partial y} = x\,(3 - x - 2y).
\]
\omterm{thm:b1:multivar:critical}{Points critiques}~: si $y = 0$, la seconde équation donne $x \in
\{0, 3\}$~; si $x = 0$, la première donne $y \in \{0, 3\}$~; si $xy
\neq 0$, on résout $2x + y = 3$, $x + 2y = 3$~: $x = y = 1$. Quatre
points~: $(0,0)$, $(3,0)$, $(0,3)$, $(1,1)$. \omterm{def:b1:derivative:def}{Dérivées} secondes~:
$r = -2y$, $s = 3 - 2x - 2y$, $t = -2x$.
\begin{itemize}
  \item $(1,1)$~: $r = -2$, $s = -1$, $t = -2$~: $rt - s^2 = 3 >
        0$, $r < 0$~: maximum local, $f(1,1) = 1$.
  \item $(0,0)$~: $r = t = 0$, $s = 3$~: $rt - s^2 = -9 < 0$~:
        \omterm{met:b1:multivar:monge}{point col}~; de même $(3, 0)$ ($s = -3$) et $(0, 3)$~: trois
        \omterm{met:b1:multivar:monge}{points cols}.
\end{itemize}
Le maximum n'est que local~: $f(-T, -T) = T^2(3 + 2T) \to
+\infty$. Contrôle par symétrie~: $f(x, y) = f(y, x)$, et de fait
l'\omterm{def:b1:logic:sets}{ensemble} \omterm{thm:b1:multivar:critical}{critique} et la classification sont symétriques en $x
\leftrightarrow y$. Interprétation~: parmi les rectangles avec
marge $x, y \geq 0$, $x + y \leq 3$, le produit $xy(3 - x - y)$ des
trois «~parts~» de $3$ est maximal quand les parts sont égales
--- une ombre à deux variables du principe de l'inégalité
arithmético-géométrique.
\end{example}

\begin{omfigure}
\begin{tikzpicture}
\begin{axis}[omaxis, width=7.6cm, height=6.6cm,
    xmin=-1.6, xmax=2.1, ymin=-1.6, ymax=2.1,
    xtick={-1,1}, ytick={-1,1},
    xlabel={$x$}, ylabel={$y$}]
  % lignes de niveau de f = x^3+y^3-3xy : f = -1 (min), f=-0.5, 0, 1
  \addplot[omDef, thick, domain=-1.35:1.9, samples=80]
    {x^2} node[pos=0.94, left, font=\small]
    {$\frac{\partial f}{\partial x} = 0$};
  \addplot[omThm, thick, domain=-1.35:1.9, samples=80]
    ({x^2}, {x}) node[pos=0.97, below right, font=\small]
    {$\frac{\partial f}{\partial y} = 0$};
  \fill (axis cs:0,0) circle (1.8pt)
    node[below left, font=\small] {col};
  \fill (axis cs:1,1) circle (1.8pt)
    node[below right, font=\small] {min};
\end{axis}
\end{tikzpicture}

{\small Les deux courbes $\nabla f = 0$ pour $f = x^3 + y^3 - 3xy$
se coupent aux \omterm{thm:b1:multivar:critical}{points critiques} $(0,0)$ (\omterm{met:b1:multivar:monge}{point col}) et $(1,1)$
(minimum local).}
\end{omfigure}

\begin{remark}[Pièges courants]
\label{rem:b1:multivar:pitfalls}
\emph{Les \omterm{def:b1:multivar:partial}{dérivées partielles} peuvent exister en un point de
discontinuité}~: le piège radial $g(x,y) = \frac{xy}{x^2+y^2}$ de
l'\cref{ex:b1:multivar:trap} vérifie $\frac{\partial g}{\partial
x}(0,0) = \frac{\partial g}{\partial y}(0,0) = 0$ (les deux
restrictions aux axes sont identiquement nulles), et pourtant $g$
n'a pas de limite à l'origine --- les \omterm{def:b1:multivar:partial}{dérivées partielles} ne
sondent que deux directions, la \omterm{def:b1:continuity:continuous}{continuité} les exige toutes~;
seule l'hypothèse $C^1$ rétablit l'ordre
(\cref{thm:b1:multivar:c1}).
\emph{Les limites selon les droites ne suffisent jamais}~: la
fonction de l'\cref{exo:b1:multivar:3} a pour limite $0$ selon toute
droite et n'a pourtant pas de limite --- toujours essayer les
paraboles (ou des majorations polaires valables uniformément en
$\theta$).
\emph{\omterm{thm:b1:multivar:critical}{Critique} est nécessaire, pas suffisant}~: les \omterm{met:b1:multivar:monge}{points cols}
abondent (trois points sur quatre dans
l'\cref{ex:b1:multivar:fourpoints})~; et le théorème ne vaut que sur
les \emph{\omterm{def:b1:topology:open}{ouverts}} --- sur un disque \omterm{def:b1:topology:closed}{fermé}, les extrema peuvent se
situer sur la \omterm{def:b1:topology:closure}{frontière} avec un \omterm{def:b1:multivar:partial}{gradient} non nul
(\cref{exo:b1:multivar:12}).
\emph{Le cas muet $rt - s^2 = 0$ est vraiment muet}~: $x^4
+ y^4$ (minimum) et $x^3 + y^3$ (ni l'un ni l'autre) ont tous deux
$r = s = t = 0$ à l'origine~; seule une étude directe du signe
tranche (\cref{exo:b1:multivar:6}, fonction $h$).
\emph{Le \omterm{def:b1:multivar:partial}{gradient} est \omterm{def:b1:euclid:orthogonal}{orthogonal} aux lignes de niveau, non porté
par elles}~: pour suivre une courbe de niveau, se déplacer
perpendiculairement à $\nabla f$~; pour monter le plus vite, se
déplacer le long de $\nabla f$ --- confondre les deux renverse la
géométrie de toute carte de niveaux.
\end{remark}

\begin{remark}[Où mènent les deux variables]
\label{rem:b1:multivar:whereused}
Ce chapitre est une porte d'entrée. Le \omterm{def:b1:multivar:partial}{gradient} et la \omterm{thm:b1:derivative:operations}{règle de la
chaîne} s'étendent mot pour mot à $n$ variables dans le volume de
Licence~2, où le $o(\norm{(h,k)})$ du \cref{thm:b1:multivar:c1}
devient la \emph{différentielle} et où le test de Monge est
démontré en entier via la \omterm{thm:b1:taylor:integral}{formule de Taylor} à l'ordre deux et les
formes quadratiques. Le cas particulier que l'on peut régler
\emph{cette} année --- les fonctions quadratiques, pour lesquelles
le développement au second ordre est exact --- fait l'objet du
devoir maison, et il se trouve que c'est le cas qui fait tourner
l'ajustement de données du monde entier~: la régression par
\omterm{pb:b1:multivar:1}{moindres carrés}. Les extrema liés (l'\cref{exo:b1:multivar:5} en
était un avant-goût) deviennent les multiplicateurs de Lagrange en
Licence~2~; les fonctions harmoniques
(\cref{exo:b1:multivar:7}) reviennent dans l'analyse complexe du
volume de Licence~3.
\end{remark}

\begin{remark}[Perspectives dans le livre 3~: l'année, bouclée]
\label{rem:b1:multivar:perspectives}
Ce chapitre est le lieu où les deux moitiés du volume se serrent
la main. La moitié analyse a fourni ses outils une \omterm{def:b1:derivative:def}{dérivée} à la
fois~: le \omterm{thm:b1:derivative:mvt}{théorème des accroissements finis} fait marcher
le \cref{thm:b1:multivar:c1}, les développements de Taylor font
marcher les tests d'extremum, et les $\varepsilon$ du
\cref{ch:b1:topology} sont revenus avec des boules au lieu
d'\omterm{prop:b1:reals:intervals}{intervalles}. La moitié algèbre a fourni la géométrie~: le
\omterm{def:b1:multivar:partial}{gradient} se lit à travers le \omterm{def:b1:euclid:def}{produit scalaire} du
\cref{ch:b1:euclid} (Cauchy--Schwarz en fait la direction de plus
forte pente), les données de Monge $(r, s, t)$ forment une matrice
symétrique du \cref{ch:b1:matrices} soumise au test du
déterminant du \cref{ch:b1:det}, et le devoir maison fait tourner
une \omterm{thm:b1:euclid:projection}{projection orthogonale} sur des vecteurs de données. Même les
courbes du \cref{ch:b1:curves} reviennent comme lignes de niveau.
Un lecteur capable de reconstituer pourquoi chacun de ces cinq
relais fonctionne a, de fait, révisé l'année entière --- ce qui
est le véritable objet de ce dernier chapitre.
\end{remark}

\begin{omfigure}
\begin{tikzpicture}[scale=1.05]
  \draw[->, thin] (-0.5,0) -- (3.7,0) node[below] {$x$};
  \draw[->, thin] (0,-0.5) -- (0,4.6) node[left] {$y$};
  \draw[omThm, very thick] (-0.28,-0.492) -- (3.3,4.52)
    node[right, font=\small] {$y = 1.4x - 0.1$};
  \fill (0,0) circle (0.06);
  \fill (1,1) circle (0.06);
  \fill (2,3) circle (0.06);
  \fill (3,4) circle (0.06);
  \draw[omDef, dashed, thick] (0,0) -- (0,-0.1);
  \draw[omDef, dashed, thick] (1,1) -- (1,1.3);
  \draw[omDef, dashed, thick] (2,3) -- (2,2.7);
  \draw[omDef, dashed, thick] (3,4) -- (3,4.1);
\end{tikzpicture}

{\small Les \omterm{pb:b1:multivar:1}{moindres carrés} en une \omterm{def:b1:linmaps:kerim}{image}~: quatre points de
données, la \omterm{pb:b1:multivar:1}{droite de régression} $y = 1.4x - 0.1$, et les résidus
verticaux (en pointillés) dont la droite minimise la somme des
carrés --- total $0.2$, le plus petit possible. Le devoir maison
calcule cette droite, démontre qu'elle est l'unique minimiseur, et
identifie toute la construction à une \omterm{thm:b1:euclid:projection}{projection orthogonale} dans
$\R^4$.}
\end{omfigure}

\section{Exercices}

\begin{exercise}[$\star$]\label{exo:b1:multivar:1}
Calculer les \omterm{def:b1:multivar:partial}{dérivées partielles}~:
$f(x,y) = x^2 y + \eu^{xy}$~; $\;g(x,y) = \ln(x^2 + y^2)$ (sur
$\R^2\setminus\{0\}$)~; $\;h(x,y) = \arctan\frac yx$ (sur $x > 0$).
\end{exercise}

\begin{exercise}[$\star$]\label{exo:b1:multivar:2}
Étudier la \omterm{def:b1:continuity:continuous}{continuité} en $(0,0)$ (où l'on pose la valeur $0$) de~:
\[
f(x,y) = \frac{x^2 y}{x^2 + y^2},
\qquad
g(x,y) = \frac{xy}{x^2 + y^2},
\qquad
h(x,y) = \frac{x^3 + y^3}{x^2 + y^2}.
\]
\emph{(Les \omterm{prop:b1:vspaces:coordinates}{coordonnées} polaires $x = \rho\cos\theta$, $y =
\rho\sin\theta$ aident~: majorer par une fonction de $\rho$ seul
quand c'est possible.)}
\end{exercise}

\begin{exercise}[$\star\star$]\label{exo:b1:multivar:3}
Soit $f(x,y) = \dfrac{x^2 y}{x^4 + y^2}$ ($f(0,0) = 0$). Démontrer
que $f$ a pour limite $0$ à l'origine \emph{selon toute droite},
mais que $f\bigl(x, x^2\bigr) = \frac12$~: $f$ n'est pas \omterm{def:b1:continuity:continuous}{continue}
en $(0,0)$.
\end{exercise}

\begin{exercise}[$\star$]\label{exo:b1:multivar:4}
Vérifier à la main le \omterm{thm:b1:multivar:schwarz}{théorème de Schwarz} sur $f(x, y) = x^3 y^2 + \sin(xy)$.
\end{exercise}

\begin{exercise}[$\star\star$]\label{exo:b1:multivar:5}
Soit $f$ de classe $C^1$ sur $\R^2$ et $g(t) = f(\cos t, \sin t)$.
Exprimer $g'(t)$ par la \omterm{thm:b1:derivative:operations}{règle de la chaîne}. En déduire que $f$
restreinte au cercle unité atteint ses extrema aux points où
$\nabla f$ est parallèle au rayon vecteur.
\end{exercise}

\begin{exercise}[$\star\star$]\label{exo:b1:multivar:6}
Trouver et classer les \omterm{thm:b1:multivar:critical}{points critiques} de~:
\[
f(x, y) = x^2 + xy + y^2 - 3x,
\qquad
g(x, y) = x^2 - y^2 + 4y,
\qquad
h(x, y) = x^4 + y^4 - 2(x - y)^2 .
\]
\emph{(Pour $h$, le test du déterminant est muet à l'origine~:
examiner $h(x, x)$ et $h(x, -x)$.)}
\end{exercise}

\begin{exercise}[$\star\star$]\label{exo:b1:multivar:7}
Une fonction $f$ est dite \emph{harmonique} lorsque
$\frac{\partial^2 f}{\partial x^2} + \frac{\partial^2 f}{\partial
y^2} = 0$. Vérifier que $x^2 - y^2$, $xy$, $\eu^x\cos y$ et
$\ln(x^2 + y^2)$ (hors de l'origine) sont harmoniques.
\end{exercise}

\begin{exercise}[$\star\star\star$]\label{exo:b1:multivar:8}
Trouver le point du plan $\{x + 2y - z = 4\} \subseteq \R^3$ le
plus proche de l'origine, de deux façons~: par \omterm{thm:b1:euclid:projection}{projection
orthogonale} (\cref{ch:b1:euclid}), et en minimisant la fonction de
deux variables $f(x, y) = x^2 + y^2 + (x + 2y - 4)^2$ obtenue en
éliminant $z$.
\end{exercise}

\begin{exercise}[$\star\star\star$]\label{exo:b1:multivar:9}
(Laplacien en \omterm{prop:b1:vspaces:coordinates}{coordonnées} polaires, premier contact) Soit $f$ de
classe $C^2$ sur $\R^2 \setminus \{0\}$ et \emph{radiale}~: $f(x, y) =
\varphi\bigl(\sqrt{x^2+y^2}\bigr)$ avec $\varphi$ de classe $C^2$
sur $\intoo{0}{+\infty}$. Démontrer que
\[
\frac{\partial^2 f}{\partial x^2} + \frac{\partial^2 f}{\partial
y^2} = \varphi''(\rho) + \frac{\varphi'(\rho)}{\rho},
\qquad \rho = \sqrt{x^2 + y^2},
\]
et trouver toutes les fonctions harmoniques radiales sur le plan
épointé.
\end{exercise}

\begin{exercise}[$\star\star$]\label{exo:b1:multivar:10}
Donner l'équation du plan tangent au graphe $z = xy$ au point
$(1, 1, 1)$. Puis trouver tous les points du graphe de
$f(x,y) = x^3 + y^3 - 3xy$ où le plan tangent est
\emph{horizontal}, et relier la réponse à
l'\cref{ex:b1:multivar:extrema}.
\end{exercise}

\begin{exercise}[$\star\star$]\label{exo:b1:multivar:11}
Soit $f$ de classe $C^1$ sur $\R^2$.
\begin{enumerate}
  \item Si $\frac{\partial f}{\partial x} = 0$ partout, démontrer
        que $f(x, y)$ ne dépend que de $y$.
  \item Trouver toutes les solutions de classe $C^1$ de l'équation
        $\frac{\partial f}{\partial x} = \frac{\partial
        f}{\partial y}$ sur $\R^2$. \emph{(Poser $g(u, v) =
        f\bigl(\frac{u+v}2, \frac{u-v}2\bigr)$ et calculer
        $\frac{\partial g}{\partial v}$ par la \omterm{thm:b1:derivative:operations}{règle de la
        chaîne}.)}
\end{enumerate}
\end{exercise}

\begin{exercise}[$\star\star\star$]\label{exo:b1:multivar:12}
Trouver le maximum et le minimum globaux de $f(x, y) = xy$ sur le
disque \omterm{def:b1:topology:closed}{fermé} $x^2 + y^2 \leq 1$. \emph{(On admettra le \omterm{thm:b1:continuity:evt}{théorème
des bornes atteintes} à deux variables~: une fonction \omterm{def:b1:continuity:continuous}{continue} sur
le disque \omterm{def:b1:topology:closed}{fermé} atteint ses bornes --- démontré dans le volume de
Licence~2. Traiter le disque \omterm{def:b1:topology:open}{ouvert} par les \omterm{thm:b1:multivar:critical}{points critiques} et le
cercle \omterm{def:b1:topology:closure}{frontière} par le paramétrage de l'\cref{exo:b1:multivar:5}.)}
\end{exercise}

\section{Problème~: moindres carrés et droite de régression}

\begin{problem}\label{pb:b1:multivar:1}
Étant donné $n$ points de données, quelle droite passe «~au plus
près~» de tous~? Legendre et Gauss ont répondu~: la droite qui
minimise la somme des \emph{carrés} des erreurs verticales ---
parce que cette minimisation-là se résout exactement par l'algèbre
linéaire. Ce problème démontre d'abord honnêtement le test de
Monge de la \cref{met:b1:multivar:monge} pour les fonctions
quadratiques (le seul cas où le développement au second ordre est
exact), puis construit les \omterm{pb:b1:multivar:1}{équations normales}, la \omterm{pb:b1:multivar:1}{droite de
régression} et le coefficient de corrélation sur la géométrie
euclidienne du
\cref{ch:b1:euclid}.\index{moindres carrés}\index{droite de régression}
\index{équations normales}\index{corrélation}

\medskip
\noindent\textbf{Partie I --- Fonctions quadratiques~: le test de
Monge, démontré.} Fixons des réels $r, s, t$ et posons $q(h, k) = r h^2 + 2s hk +
t k^2$.
\begin{enumerate}
  \item Supposons $r \neq 0$. Établir la forme canonique
        \[
        q(h, k) = r\Bigl(h + \frac{s}{r}k\Bigr)^{\!2}
        + \frac{rt - s^2}{r}\,k^2 ,
        \]
        et en déduire~: si $rt - s^2 > 0$, $q$ a le signe strict
        de $r$ hors de l'origine~; si $rt - s^2 < 0$, $q$ prend
        les deux signes.
  \item Régler les cas restants~: $r = 0$, $t \neq 0$
        (symétrique)~; et $r = t = 0$ ($q = 2shk$). Conclure~:
        $q$ prend les deux signes si et seulement si $rt - s^2 <
        0$, et $q$ ne s'annule qu'à l'origine si et seulement si
        $rt - s^2 > 0$.
  \item Soit maintenant $f(x, y) = \frac12\bigl(r x^2 + 2s xy + t
        y^2\bigr) + \beta x + \gamma y + \delta$ une fonction
        quadratique ayant un \omterm{thm:b1:multivar:critical}{point critique} $X_0 = (a, b)$.
        Démontrer le développement \emph{exact}
        \[
        f(X_0 + (h,k)) = f(X_0) + \tfrac12\,q(h, k)
        \qquad (\text{sans reste}),
        \]
        et en déduire la classification de Monge pour les
        fonctions quadratiques~: minimum global strict si $rt -
        s^2 > 0$, $r > 0$~; maximum global strict si $rt - s^2 >
        0$, $r < 0$~; \omterm{met:b1:multivar:monge}{point col} si $rt - s^2 < 0$.
  \item Supposons $rt - s^2 > 0$ et $r > 0$ (donc aussi $t > 0$).
        Démontrer la minoration de coercivité explicite
        \[
        q(h, k) \geq c\,(h^2 + k^2),
        \qquad
        c = \frac{rt - s^2}{r + t} > 0 .
        \]
        \emph{(Montrer que la forme quadratique $q - c(h^2 + k^2)$,
        de coefficients $r - c$, $s$, $t - c$, satisfait encore
        aux conditions du test du discriminant.)}
  \item En déduire qu'une fonction quadratique dont la partie
        quadratique est définie positive tend vers $+\infty$
        quand $\norm{(x,y)} \to \infty$, et admet donc un
        \emph{unique} minimiseur \emph{global}~: son \omterm{thm:b1:multivar:critical}{point
        critique}. Contraster avec
        l'\cref{ex:b1:multivar:extrema}, où un minimum local d'une
        fonction non quadratique n'était pas global.
\end{enumerate}

\medskip
\noindent\textbf{Partie II --- Les \omterm{pb:b1:multivar:1}{équations normales}.} Soient
$C_1, C_2, b \in \R^n$ (\omterm{def:b1:euclid:def}{produit scalaire} canonique), et
\[
f(u, v) = \norm{u\,C_1 + v\,C_2 - b}^2 .
\]
\begin{enumerate}[resume]
  \item Développer $f$ et montrer que c'est une fonction
        quadratique de $(u, v)$ dont la partie quadratique a pour
        coefficients $r = 2\norm{C_1}^2$,
        $s = 2\langle C_1, C_2\rangle$, $t = 2\norm{C_2}^2$.
  \item Montrer que $rt - s^2 > 0$ si et seulement si $(C_1,
        C_2)$ est \omterm{def:b1:vspaces:free}{libre} \emph{(le cas d'égalité de
        Cauchy--Schwarz, \cref{thm:b1:euclid:cs})}~; on le
        supposera désormais.
  \item Montrer que les équations \omterm{thm:b1:multivar:critical}{critiques} $\nabla f = 0$ sont
        les \emph{\omterm{pb:b1:multivar:1}{équations normales}}
        \[
        \begin{pmatrix}
        \norm{C_1}^2 & \langle C_1, C_2\rangle\\
        \langle C_1, C_2\rangle & \norm{C_2}^2
        \end{pmatrix}
        \begin{pmatrix} u\\ v\end{pmatrix}
        =
        \begin{pmatrix} \langle C_1, b\rangle\\ \langle C_2,
        b\rangle\end{pmatrix}
        \]
        --- la matrice de Gram de l'\cref{exo:b1:euclid:11} à
        gauche --- et qu'elles disent exactement~: $b - (uC_1 +
        vC_2) \perp C_1, C_2$. Conclure avec la partie I et
        le \cref{thm:b1:euclid:projection}~: l'unique minimiseur
        donne $p = uC_1 + vC_2 = $ la \omterm{thm:b1:euclid:projection}{projection orthogonale} de
        $b$ sur $\operatorname{Vect}(C_1, C_2)$.
  \item Montrer que la valeur minimale est $\norm{b}^2 - \langle
        p, b\rangle = \norm{b - p}^2$, et faire le dessin
        pythagoricien.
\end{enumerate}

\medskip
\noindent\textbf{Partie III --- La \omterm{pb:b1:multivar:1}{droite de régression}.} Points
de données $(x_1, y_1), \dots, (x_n, y_n)$, les $x_i$ n'étant pas
tous égaux. Minimiser
\[
E(\alpha, \beta) = \sum_{i=1}^{n}\bigl(y_i - \alpha - \beta
x_i\bigr)^2 .
\]
On note $\bar x = \frac1n\sum x_i$, $\bar y = \frac1n\sum y_i$,
$v_x = \frac1n\sum x_i^2 - \bar x^2$, $v_y = \frac1n\sum y_i^2 -
\bar y^2$, $c_{xy} = \frac1n\sum x_i y_i - \bar x\bar y$.
\begin{enumerate}[resume]
  \item Reconnaître la partie II avec $C_1 = (1, \dots, 1)$, $C_2
        = (x_1, \dots, x_n)$, $b = (y_1, \dots, y_n)$, et écrire
        les \omterm{pb:b1:multivar:1}{équations normales}
        \[
        n\alpha + \Bigl(\sum x_i\Bigr)\beta = \sum y_i,
        \qquad
        \Bigl(\sum x_i\Bigr)\alpha + \Bigl(\sum
        x_i^2\Bigr)\beta = \sum x_i y_i .
        \]
  \item Les résoudre~:
        \[
        \beta = \frac{c_{xy}}{v_x},
        \qquad
        \alpha = \bar y - \beta\,\bar x ,
        \]
        et observer que la \emph{\omterm{pb:b1:multivar:1}{droite de régression}} $y = \alpha
        + \beta x$ passe par le point moyen $(\bar x, \bar
        y)$.
  \item Vérifier que $v_x > 0$ exactement quand les $x_i$ ne sont
        pas tous égaux, et rapprocher cela de la question 7.
  \item Démontrer que l'erreur minimale vaut
        \[
        E_{\min} = n\Bigl(v_y - \frac{c_{xy}^2}{v_x}\Bigr)
        = n\,v_y\,(1 - \rho^2),
        \qquad
        \rho = \frac{c_{xy}}{\sqrt{v_x v_y}}
        \quad (v_y \neq 0).
        \]
        En déduire $\abs\rho \leq 1$, avec $\abs\rho = 1$ si et
        seulement si les données sont parfaitement alignées.
  \item Exemple traité~: pour les données $(0,0)$, $(1,1)$,
        $(2,3)$, $(3,4)$, calculer $\bar x, \bar y, v_x, c_{xy}$,
        la \omterm{pb:b1:multivar:1}{droite de régression}, les quatre résidus et $E_{\min}$.
  \item Démontrer que les résidus $\varepsilon_i = y_i - \alpha -
        \beta x_i$ de la droite optimale vérifient $\sum_i
        \varepsilon_i = 0$ et $\sum_i x_i\varepsilon_i = 0$,
        et interpréter les deux par l'orthogonalité.
\end{enumerate}

\medskip
\noindent\textbf{Partie IV --- Variantes et \omterm{def:b1:logic:map}{applications}.}
\begin{enumerate}[resume]
  \item (Meilleure constante) Montrer que la constante $\alpha$
        qui minimise $\sum_i (y_i - \alpha)^2$ est la moyenne
        $\bar y$, et que la valeur minimale est $n\,v_y$~: la
        variance mesure le défaut de constance des données.
  \item (Droite passant par l'origine) Montrer que la pente qui
        minimise $\sum_i (y_i - \beta x_i)^2$ est $\beta_0 =
        \frac{\sum x_iy_i}{\sum x_i^2}$, et donner une condition
        sur les données pour que $\beta_0$ coïncide avec la pente
        $c_{xy}/v_x$ de la question 11.
  \item (Distance à une droite, encore) Pour un point $M$ et la
        droite $D$ passant par $P_0$ et dirigée par le vecteur
        unitaire $w$, minimiser $g(\tau) = \norm{P_0 + \tau w -
        M}^2$ et en déduire $d(M, D)^2 = \norm{M - P_0}^2 -
        \langle M - P_0, w\rangle^2$~; retrouver la formule $d =
        \frac{\abs{a x_0 + b y_0 + c}}{\sqrt{a^2 + b^2}}$ pour la
        droite $ax + by + c = 0$ du plan.
  \item (Analyse de la variance) Démontrer la décomposition $v_y =
        \beta^2 v_x + \frac{E_{\min}}n$~: la variance des $y_i$ se
        scinde en la part expliquée par la droite et la variance
        résiduelle.
  \item (Ajustement parabolique) Pour ajuster $y = a + bx + cx^2$,
        montrer que les \omterm{pb:b1:multivar:1}{équations normales} forment le système $3
        \times 3$ de matrice des moments
        \[
        \begin{pmatrix}
        n & \sum x_i & \sum x_i^2\\
        \sum x_i & \sum x_i^2 & \sum x_i^3\\
        \sum x_i^2 & \sum x_i^3 & \sum x_i^4
        \end{pmatrix},
        \]
        matrice de Gram inversible dès que trois des $x_i$ sont
        distincts \emph{(liberté de $(1, X, X^2)$ échantillonné
        aux données, \cref{exo:b1:euclid:11}~; comparer avec les
        matrices des moments du devoir maison du
        \cref{ch:b1:det})}.
\end{enumerate}

\medskip
\noindent\textbf{Partie V --- Robustesse, et synthèse.}
\begin{enumerate}[resume]
  \item Classer les \omterm{thm:b1:multivar:critical}{points critiques} des trois fonctions
        quadratiques
        \[
        x^2 + xy + y^2, \qquad x^2 + 3xy + y^2, \qquad
        x^2 + 2xy + y^2 ,
        \]
        par la partie I, en traitant à la main le troisième cas
        dégénéré (où le minimum est-il atteint~?).
  \item (Expérience de la valeur aberrante) Ajouter le point $(10,
        0)$ aux données de la question 14 et recalculer la pente
        $\beta$. Que s'est-il passé, et pourquoi l'erreur
        \emph{quadratique} est-elle si sensible à un seul point
        lointain~?
  \item (\omterm{pb:b1:multivar:1}{Moindres carrés} pondérés) Étant donné des poids $w_i >
        0$, minimiser $\sum_i w_i(y_i - \alpha - \beta x_i)^2$~:
        montrer que la solution est donnée par les mêmes formules
        qu'à la question 11 avec des moyennes, une variance et une
        covariance pondérées (les définir).
  \item Montrer que, pour la droite optimale, $E_{\min} = 0$
        force tous les points à être exactement sur elle, et
        relier au cas d'égalité $\abs\rho = 1$ de la question
        13~; tester sur les données alignées $(1,1), (2,3),
        (3,5)$.
  \item Synthèse, en quatre phrases~: pourquoi les fonctions
        quadratiques sont la seule classe où le test du second
        ordre de ce chapitre n'a besoin d'aucun théorème admis~;
        comment les \omterm{pb:b1:multivar:1}{équations normales} identifient la
        minimisation analytique à la \omterm{thm:b1:euclid:projection}{projection orthogonale} du
        \cref{ch:b1:euclid}~; ce que mesure le coefficient de
        corrélation et quelle inégalité le borne~; et quels
        morceaux des chapitres 18 à 23 (\omterm{def:b1:vspaces:free}{bases}, matrices de Gram
        et des moments, \omterm{def:b1:linmaps:projection}{projections}) ont reparu. Nommer la
        méthode et les équations étudiées dans les parties II et
        III.
\end{enumerate}
\end{problem}
